\documentclass{situationist-preprint}

\title{On Imaginative Resistance}
\msid{AES-08-b}
\msfamily{AES-08}
\msdiscipline{Aesthetics}
\mstype{Companion manuscript}
\mskeywords{imaginative resistance; moderate moralism; prescribed response; fictional truth; Hume}

\begin{document}
\maketitle

\begin{abstract}
Readers who imagine dragons and time travel without difficulty balk at imagining that, in a story, a gratuitous killing was right. This is the phenomenon of imaginative resistance. AES-08-a holds that a moral flaw in a work is an aesthetic flaw when, and only when, it prevents a sensitive audience from giving a response the work prescribes. We argue that imaginative resistance comes in two forms, resistance to a fictional truth and resistance to a prescribed response, and that only the second is the phenomenon on which moderate moralism turns. The distinction explains why some instances of imaginative resistance mark aesthetic flaws and others do not, and it locates Hume's observation about vicious manners on the correct side.
\end{abstract}

\section{The puzzle}

Hume observed that he could not enter into sentiments of a work that approved of vicious manners: where a poet's representation of manners is not marked with disapproval, the work is deformed, and he will not relish the composition on that account \citep{Hume1757}. \citet{Moran1994}, \citet{Walton1994} and \citet{Gendler2000} revived the observation as a puzzle. Why can a reader imagine that in a story Hector was killed by a god, but not that in a story it was right for Giselda to kill her baby because it was a girl, when the author asserts as much?

Two answers dominate. On the first, the resistance is a matter of inability: moral facts supervene on non-moral facts, and a story cannot make a moral claim fictionally true without making true the non-moral facts that would ground it \citep{Walton1994}. On the second, the resistance is a matter of unwillingness: readers decline to imagine the claim because doing so would amount to adopting, as their own, a moral perspective they reject \citep{Gendler2000}.

\section{Two forms of resistance}

\begin{definition}\label{d:forms}
\emph{Propositional resistance} is reluctance or inability to accept that a moral proposition the work asserts is true in the fiction. \emph{Responsive resistance} is reluctance or inability to give a response the work prescribes, in the sense of Definition~2.1 of AES-08-a, because the response would require endorsing a moral perspective the reader recognizes as unacceptable.
\end{definition}

The two often occur together but are distinct. A reader may accept that, in a story told by a monstrous narrator, the narrator judges the killing right, and resist nothing, since the work prescribes horror at the narrator. Conversely, a reader may grant without difficulty that in the fiction a character is admirable, yet be unable to admire him, because the admiration the work prescribes depends on finding his contempt for others attractive.

\begin{proposition}\label{p:resp}
Responsive resistance by a sensitive audience is an aesthetic flaw in the work; propositional resistance alone is not.
\end{proposition}

\begin{proof}
Responsive resistance by a sensitive audience is, by Principle~P2 of AES-08-a, the failure of such an audience to give a prescribed response because it would require endorsing a perspective recognized as morally unacceptable. By Principle~P1, a work fails as the kind of work it is to the extent that a sensitive and informed audience cannot give the responses it prescribes, and by Theorem~3.1 a moral flaw that prevents such a response is an aesthetic flaw. Propositional resistance concerns what is true in the fiction. A reader who cannot accept that a moral claim is fictionally true may still give every response the work prescribes, for instance where the work prescribes contempt for the character who asserts the claim. In that case nothing the work prescribes is blocked, and by Theorem~3.1 any moral flaw involved is a moral flaw only.
\end{proof}

\begin{corollary}\label{c:hume}
Hume's deformity is responsive resistance.
\end{corollary}

\begin{proof}
What Hume describes is not an inability to suppose that the manners depicted were approved of in the world of the poem, which he grants, but an inability to enter into the sentiments the poem asks its reader to share. That is a prescribed response that a sensitive reader cannot give, and by Proposition~\ref{p:resp} it is an aesthetic flaw.
\end{proof}

\begin{objection}[Inability or refusal]
Moderate moralism needs the audience not to be \emph{able} to give the response; Gendler's account makes resistance a matter of will. If the reader merely refuses, the work has not failed; the reader has.
\end{objection}

\begin{reply}
Principle~P2 of AES-08-a is a claim about what a sensitive audience does, not about what it could do by an effort of will. A sensitive reader could perhaps feign amusement or work herself into admiration, as one can feign laughter at a cruel joke. But the response the work prescribes is one that is merited by the work (Definition~2.2 of AES-08-a), and a response produced by an act of will against one's moral perception is not the merited response the work aims at. Whether the resistance is called inability or refusal, the work does not secure the response it prescribes from the audience it is addressed to, and that is a failure on its own terms.
\end{reply}

\section{Conclusion}

Imaginative resistance is not one phenomenon. Resistance to what is fictionally true is a puzzle about fiction and may carry no aesthetic consequence. Resistance to what a work asks us to feel is the mechanism by which, on the moderate moralism of AES-08-a, a moral flaw becomes an aesthetic one.

\begin{thebibliography}{99}\small
\bibitem[Gendler(2000)]{Gendler2000} Gendler, T.~S. (2000). The puzzle of imaginative resistance. \emph{The Journal of Philosophy}, 97(2), 55--81.
\bibitem[Hume(1757)]{Hume1757} Hume, D. (1757). Of the standard of taste. In \emph{Four Dissertations}. A. Millar.
\bibitem[Moran(1994)]{Moran1994} Moran, R. (1994). The expression of feeling in imagination. \emph{The Philosophical Review}, 103(1), 75--106.
\bibitem[Walton(1994)]{Walton1994} Walton, K.~L. (1994). Morals in fiction and fictional morality. \emph{Proceedings of the Aristotelian Society, Supplementary Volume}, 68, 27--50.
\end{thebibliography}

\end{document}
